\chapter{黎曼度量}

\section{黎曼度量}

令 $M$ 是光滑流形。$M$ 上的\emph{黎曼度量}指的是一个
光滑协变 $2$-张量场 $g\in \mathcal{T}^2(M)$，其在每个点 $p\in M$
处都是 $T_pM$ 上的一个内积。因此，$g$ 是一个对称 $2$-张量场，
并且是正定的，即对于每个 $p\in M$ 和 $v\in T_pM$，有 $g_p(v,v)\geq 0$
，且当且仅当 $v=0$ 时取等号。一个\emph{黎曼流形}指的是
一个配备了黎曼度量的光滑流形 $(M,g)$。如果 $M$ 上的黎曼度量
是清楚的，通常也说 $M$ 是一个黎曼流形。

下面的命题表明黎曼度量是广泛存在的，在 Lee 的《光滑流形导引》中
有具体的证明。

\begin{proposition}
  每个光滑流形都有一个黎曼度量。
\end{proposition}


令 $g$ 是带边或者无边光滑流形 $M$ 上的一个黎曼度量。对于每个 $p\in M$
和 $v,w\in T_pM$，我们通常采用记号
\[
  \inn{v,w}_g=g_p(v,w).
\]
利用内积，我们可以定义向量 $v\in T_pM$ 的长度为 $|v|_g=\sqrt{\inn{v,v}_g}$。
如果度量是清楚的，我们也可以省略下标 $g$。

\begin{example}[Euclid 度量]
  $\mathbb{R}^n$ 上的\emph{Euclid 度量} $\bar g$ 就是通常意义下的点积，
  对于每个 $x\in \mathbb{R}^n$，切空间 $T_x \mathbb{R}^n$ 可以自然地
  同构为 $\mathbb{R}^n$ 本身。这意味着对于 $v,w\in T_x \mathbb{R}^n$，
  在标准坐标 $(x^1,\dots,x^n)$ 中有 $v=\sum v^i\partial_i|_x$ 和
  $w=\sum w^i\partial_i|_x$，因此
  \[
    \inn{v,w}_{\bar g}= \sum_{i=1}^n v^i w^i.
  \]
  以后，当谈及 $\mathbb{R}^n$ 时，我们总是指的是 $\mathbb{R}^n$ 配备 Euclid 度量。
\end{example}

\section{等距}

设 $(M,g)$ 和 $\bigl(\wtilde M,\tilde g\bigr)$ 是带边或者无边黎曼流形。
$(M,g)$ 到 $\bigl(\wtilde M,\tilde g\bigr)$ 的\emph{等距}指的是
一个微分同胚 $\varphi:M\to\wtilde M$，并且使得 $\varphi^*\tilde g=g$。
这意味着每个微分 $d\varphi_p:T_pM\to T_{\varphi(p)}\wtilde M$
都是线性等距。如果存在这样一个等距映射，那么我们就说 $(M,g)$ 和
$\bigl(\wtilde M,\tilde g\bigr)$ 是\emph{等距的}。

显然等距的复合和逆也都是等距，所以等距构成一个等价关系。黎曼几何
的课题就是考虑黎曼流形的哪些性质被等距所保持。

如果 $(M,g)$ 和 $\bigl(\wtilde M,\tilde g\bigr)$ 是黎曼流形，每个点 $p\in M$
都有一个邻域 $U$ 使得 $\varphi|_U$ 是到 $\wtilde M$ 的某个开子集的等距，那么
我们说 $\varphi:M\to\wtilde M$ 是\emph{局部等距}。

一个黎曼 $n$-流形被称为\emph{平坦的}，如果其局部等距于某个 Euclid 空间，也即
每个点都有一个邻域和 $\mathbb{R}^n$ 中的某个开子集等距。可以证明所有黎曼
$1$-流形都是平坦的。

$(M,g)$ 到自身的等距被称为 $(M,g)$ 的等距，显然所有 $(M,g)$ 的等距在复合下
构成一个群，被称为 $(M,g)$ 的\emph{等距群}，记为 $\Iso(M,g)$。有时候
也简记为 $\Iso(M)$。

\section{度量的局部表示}

令 $(M,g)$ 是带边或者无边黎曼流形。$(x^1,\dots,x^n)$ 是开子集 $U\subseteq M$
上的光滑局部坐标，那么 $g$ 可以局部表示为
\[
  g=g_{ij}dx^i\otimes dx^j.
\]
其中 $g_{ij}:U\to \mathbb{R}$ 是光滑函数。$g$ 也可以用对称积表示，利用 $g_{ij}$
的对称性，有
\begin{align*}
  g&=g_{ij}dx^i\otimes dx^j =\frac{1}{2}(g_{ij}dx^i\otimes dx^j+g_{ji}dx^i\otimes dx^j)\\
  &=\frac{1}{2}g_{ij}(dx^i\otimes dx^j+dx^j\otimes dx^i)=g_{ij}dx^i dx^j.
\end{align*}

关于 $M$ 的在开集 $U$ 上的局部标架 $(E_i)$ 被称为\emph{正交标架}，如果在每个 $p\in U$
处向量 $E_1|_p,\dots,E_n|_p$ 构成 $T_pM$ 的一个正交基。等价地说，$(E_{i})$
是正交标架当且仅当
\[
  \inn{E_i,E_j}=\delta_{ij},
\]
此时 $g$ 可以局部表示为
\[
  g=(\varepsilon^1)^2+\cdots+(\varepsilon^n)^2,
\]
其中 $(\varepsilon^i)^2$ 表示对称积 $\varepsilon^i\varepsilon^i=\varepsilon^i\otimes\varepsilon^i$。

\begin{proposition}[正交标架的存在性]
  令 $(M,g)$ 是带边或者无边黎曼 $n$-流形。如果 $(X_j)$ 是开集 $U\subseteq M$ 上的光滑局部标架，
  那么存在 $U$ 上的光滑正交标架 $(E_j)$ 使得：对于每个 $k=1,\dots,n$ 和 $p\in U$，
  有 $\mathrm{span}(E_1|_p,\dots,E_k|_p)=\mathrm{span}(X_1|_p,\dots,E_k|_p)$。
  特别地，这表明对于每个 $p\in M$，都存在一个定义在 $p$ 的某个邻域上的局部正交标架。
\end{proposition}
\begin{proof}
  对于每个 $p\in U$，对向量组 $(X_1|_p,\dots,X_n|_p)$ 应用 Gram-Schmidt 正交化，
  得到 $U$ 上的正交向量场 $(E_1,\dots,E_n)$，且满足上面的张成条件。
  由于 Gram-Schmidt 正交化过程是光滑的（因为 $g$ 是光滑的），
  得到的正交向量场 $(E_1,\dots,E_n)$ 也是光滑的。
  因此，$(E_j)$ 是满足命题条件的光滑正交标架。
\end{proof}

注意，这并不意味着我们可以找到 $p$ 附近的某个坐标使得坐标标架 $(\partial_i)$
是正交的。实际上，这只在平坦度量的情况下成立。

对于带边或者无边黎曼流形 $(M,g)$，我们可以定义\emph{单位切丛} $UTM\subseteq TM$
为由单位向量构成的子集：
\begin{equation}
  UTM=\{(p,v)\in TM\,|\, |v|_g=1\}.
\end{equation}

\begin{proposition}
如果 $(M,g)$ 是带边或者无边黎曼流形，那么单位切丛 $UTM$ 是光滑的，并且
是 $TM$ 的恰当嵌入的余维数为 $1$ 的带边子流形，边界 $\partial(UTM)=\pi^{-1}(\partial M)$，
这里 $\pi:UTM\to M$ 是自然投影。单位切丛连通当且仅当 $M$ 连通，单位切丛紧当且仅当 $M$ 紧。
\end{proposition}
\begin{proof}
  定义映射 $F:TM\to \mathbb{R}$ 为 $F(p,v)=|v|_g^2$。任取光滑坐标卡 $(x^i)$，
  此时 $v=v^i\partial_i|_p$，因此 $F$ 有局部表示
  \[
    (x^1,\dots,x^n,v^1,\dots,v^n)\mapsto g_{ij}(x)v^iv^j,
  \]
  这是光滑的，所以 $F$ 是光滑映射。显然 $F$ 是浸没，注意到 $UTM=F^{-1}(1)$，所以 $UTM$ 是 $TM$ 的
  余维数为 $1$ 的嵌入子流形。任取 $p\in\partial M$，则 $T_pM$ 是带边流形 $M$ 的切空间，所以 $T_pM$ 是带边空间，其边界为
  $T_p(\partial M)$。因此，$UTM\cap T_pM$ 是 $T_pM$ 的单位球面与边界 $T_p(\partial M)$ 的交集，所以
  $\partial(UTM)=\pi^{-1}(\partial M)$。由于每个 $T_pM$ 都是连通的，所以 $UTM$ 连通当且仅当 $M$ 连通。由于 $\pi:UTM\to M$ 是连续的且纤维是紧的，所以 $UTM$ 紧当且仅当 $M$ 紧。
\end{proof}



\section{构造黎曼流形的方法}

\subsection{黎曼子流形}

黎曼流形的每个子流形都自动继承一个黎曼度量，这来源于下面的引理。

\begin{lemma}
  设 $\bigl(\wtilde M,\tilde g\bigr)$ 是带边或者无边黎曼流形，
  $M$ 是带边或者无边光滑流形，$F:M\to\wtilde M$ 是光滑映射。
  $2$-张量场 $g=F^*\tilde{g}$ 是 $M$ 上的一个黎曼度量当且仅当
  $F$ 是浸入。
\end{lemma}
\begin{proof}
  若 $g=F^*\tilde{g}$ 是黎曼度量，任取 $p\in M$，如果 $dF_p(v)=0$，那么
  \[
    \inn{v,v}_g=\tilde g_{F(p)}(dF_p(v),dF_p(v))=0,
  \] 
  所以 $v=0$，因此 $dF_p$ 是单射，所以 $F$ 是浸入。
  反之，若 $F$ 是浸入，任取 $p\in M$ 和 $v\in T_pM$，如果 $v\neq 0$，那么
  $dF_p(v)\neq 0$，所以
  \[
    \inn{v,v}_g=\tilde g_{F(p)}(dF_p(v),dF_p(v))>0,
  \]
  因此 $g$ 是正定的，所以 $g$ 是黎曼度量。
\end{proof}

假设 $\bigl(\wtilde M,\tilde g\bigr)$ 是带边或者无边黎曼流形。
给定光滑浸入 $F:M\to\wtilde M$，度量 $g=F^*\tilde g$ 被称为\emph{$F$ 的诱导度量}。
另一方面，如果 $M$ 已经配备了一个黎曼度量 $g$，那么满足 $F^*\tilde g=g$ 的
浸入或者嵌入 $F:M\to\wtilde M$ 被称为\emph{等距浸入}或者\emph{等距嵌入}。
使用哪个术语取决于 $M$ 上的度量是否已经被指定。

诱导度量最重要的例子是子流形。假设 $M\subseteq \wtilde M$
是(浸入或者嵌入)子流形。\emph{$M$ 上的诱导度量}指的是由包含映射
$\iota:M\hookrightarrow\wtilde M$ 诱导的度量 $g=\iota^*\tilde g$。
在这个度量下，$M$ 被称为 $\wtilde M$ 的\emph{黎曼子流形}。

如果 $(M,g)$ 是 $\bigl(\wtilde M,\tilde g\bigr)$ 的黎曼子流形，
那么对于每个 $p\in M$ 和 $v,w\in T_pM$，有
\[
  g_p(v,w)=\tilde g_p\bigl(d\iota_p(v),d\iota_p(w)\bigr).
\]
我们通常把 $T_pM$ 视为 $T_p\wtilde M$ 的子空间 $d\iota_p(T_pM)$，
因此上式通常写为 $g_p(v,w)=\tilde g_p(v,w)$。换句话说，诱导度量
$g$ 仅仅是 $\tilde g$ 在 $M$ 上的限制。

\begin{example}[球面]
  对于每个正整数 $n$，单位 $n$-球面 $\mathbb{S}^n\subseteq \mathbb{R}^{n+1}$
  是嵌入子流形。由 Euclid 度量诱导的 $\mathbb{S}^n$ 上的黎曼度量记为
  $\mathring g$，被称为 $\mathbb{S}^n$ 上的\emph{圆度量}或者说\emph{标准度量}。
\end{example}

如果 $\bigl(\wtilde M,\tilde g\bigr)$ 是 $m$-维黎曼流形，$M\subseteq \wtilde M$
是 $n$-维子流形，开子集 $\wtilde U\subseteq \wtilde M$ 上的一个局部标架
$(E_1,\dots,E_m)$ 被称为\emph{适应 $M$ 的}，如果前 $n$ 个向量场 $(E_1,\dots,E_n)$ 
与 $M$ 相切。在 $\wtilde M$ 无边的情况下，适应 $M$ 的局部标架总是存在的。

\begin{proposition}[适应正交标架的存在性]\label{prop:adapted orthonormal frame}
  令 $\bigl(\wtilde M,\tilde g\bigr)$ 是无边黎曼流形，$M\subseteq \wtilde M$ 是
  嵌入的带边或者无边子流形。给定 $p\in M$，存在 $p$ 在 $\wtilde M$ 中的
  某个邻域 $\wtilde U$ 和 $\wtilde U$ 上的适应 $M$ 的光滑正交标架。
\end{proposition}
\begin{proof}
  设 $\wtilde M$ 是 $m$ 维，$M$ 是 $n$ 维子流形。任取 $p\in M$，
  取 $p$ 在 $\wtilde M$ 中的切片坐标卡 $(\wtilde U,(x^1,\dots,x^m))$，使得
  $M\cap \wtilde U$ 由方程 $x^{n+1}=\cdots=x^m=0$ 给出。对这个坐标标架
  $(\partial_i)$ 应用 Gram-Schmidt 正交化，得到 $\wtilde U$ 上的光滑正交标架
  $(E_1,\dots,E_m)$。由于正交化过程保持张成性质，所以 $(E_1,\dots,E_n)$ 与 $M$ 相切。
  于是 $(E_1,\dots,E_m)$ 是适应 $M$ 的光滑正交标架。
\end{proof}

假设 $\bigl(\wtilde M,\tilde g\bigr)$ 是黎曼流形，$M\subseteq \wtilde M$ 是
带边或者无边子流形。给定 $p\in M$，向量 $v\in T_p\wtilde M$ 被称为\emph{正交于 $M$ 的}，
如果对于每个 $w\in T_pM$，有 $\inn{v,w}=0$。$T_p\wtilde M$ 的所有正交于 $M$ 的向量
构成 $T_p\wtilde M$ 的一个子空间，记为 $N_pM=(T_pM)^\bot$，被称为\emph{$p$ 处的法空间}。
在每个 $p\in M$ 处，环境切丛 $T_p\wtilde M$ 可以正交分解为
$T_p\wtilde M=T_pM\oplus N_pM$。环境切丛 $T\wtilde M|_M$ 的截面
$N$ 被称为\emph{沿 $M$ 的法向量场}，如果对于每个 $p\in M$，有 $N_p\in N_pM$。
集合
\[
  NM=\coprod_{p\in M}N_pM
\]
被称为 $M$ 的\emph{法丛}。

\begin{proposition}[法丛]\label{prop:normal bundle}
  如果 $\wtilde M$ 是黎曼 $m$-流形，$M\subseteq \wtilde M$ 是浸入或者嵌入的
  黎曼 $n$-子流形，那么 $NM$ 是环境切丛 $T\wtilde M|_M$ 的
  秩 $(m-n)$ 的光滑子丛。光滑丛同态
  \[
    \pi^\top:T\wtilde M|_M\to TM,\quad \pi^\bot:T\wtilde M|_M\to NM
  \]
  分别被称为 \emph{切向投影}和\emph{法向投影}，它们在每个点 $p\in M$ 处
  分别限制为 $T_p\wtilde M$ 到 $T_pM$ 和 $N_pM$ 的正交投影。
\end{proposition}
\begin{proof}
  给定 $p\in M$，那么存在 $p$ 处的邻域 $U\subseteq M$ 使得 $U\subseteq \wtilde M$
  是嵌入子流形。根据 \autoref{prop:adapted orthonormal frame}，在 $p$
  的某个邻域 $\wtilde U\subseteq \wtilde M$ 上存在适应 $U$ 的光滑正交标架
  $(E_1,\dots,E_m)$。于是 $(E_1,\dots,E_n)$ 是 $\wtilde U\cap U$
  上的正交标架。那么后 $m-n$ 个向量场 $(E_{n+1},\dots,E_m)$ 
  构成了 $NM$ 的光滑局部标架，这就表明 $NM$ 是光滑子丛。

  丛同态 $\pi^\top$ 和 $\pi^\bot$ 在每个点 $p\in M$ 处的定义就是向
  切空间和法空间做正交投影，所以它们是唯一确定的。用适应的正交标架，它们可以
  表示为
  \begin{align*}
    \pi^\top(X^1E_1+\cdots+X^mE_m)&=X^1E_1+\cdots+X^nE_n,\\
    \pi^\bot(X^1E_1+\cdots+X^mE_m)&=X^{n+1}E_{n+1}+\cdots+X^mE_m.
  \end{align*}
  这表明它们都是光滑的。 
\end{proof}

在 $\wtilde M$ 是带边流形的情况下，上面的构造不总是成立，
因为此时没有切片坐标卡的一般构造。但是，当子流形是边界本身时，
利用边界条件而不是切片坐标卡，我们仍然可以构造法丛。

假设 $(M,g)$ 是带边黎曼流形。我们总是把边界 $\partial M$ 视为 $M$ 的
黎曼子流形，附带诱导度量。

\begin{proposition}[向外法向的存在性]
  如果 $(M,g)$ 是带边黎曼流形，那么 $\partial M$ 的法丛是
  秩 $1$ 的光滑向量丛，并且存在唯一沿 $\partial M$ 的向外的单位光滑法向量场。
\end{proposition}

在子流形 $M\subseteq \wtilde M$ 上计算通常使用\emph{光滑局部参数化}：
一个光滑映射 $X:U\to\wtilde M$，其中 $U$ 是 $\mathbb{R}^n$
的开子集，使得 $X(U)$ 是 $M$ 的开子集，并且 $X$ 视为 $U\to M$
的映射时是到像集的微分同胚。注意到我们可以把 $X$ 视为到 $M$ 或者
到 $\wtilde M$ 的映射，我们都记为 $X$。显然，如果记 $V=X(U)$，
$\varphi=X^{-1}:V\to U$，那么 $(V,\varphi)$ 是 $M$ 的一个光滑坐标卡。

假设 $(M,g)$ 是 $\bigl(\wtilde M,\tilde g\bigr)$ 的黎曼子流形，
$X:U\to\wtilde M$ 是 $M$ 的一个光滑局部参数化。$g$ 在这个坐标中
的局部表示为 $U$ 上的 $2$-张量场：
\[
  (\varphi^{-1})^*g=X^*g=X^*\iota^*\tilde g=(\iota\circ X)^*\tilde g.
\]
因为 $\iota\circ X$ 就是 $X$ 本身，所以这实际上就是 $X^*\tilde g$。
然后只需要应用拉回张量的计算公式即可。例如，如果 $M$ 是 $\mathbb{R}^m$
的浸入 $n$-维黎曼子流形，$X:U\to \mathbb{R}^m$ 是 $M$ 的光滑局部参数化，
那么 $U$ 上的诱导度量为 
\[
  g=X^*\bar g=\sum_{i=1}^m (dX^i)^2
  =\sum_{i=1}^m\left(
    \sum_{j=1}^m\frac{\partial X^i}{\partial u^j}du^j
  \right)^2.
\]

\begin{example}[图像坐标中的度量]
  如果 $U\subseteq \mathbb{R}^n$ 是开集，$f:U\to \mathbb{R}$ 是光滑函数，
  那么 $f$ 的图像指的是子集 $\Gamma(f)=\{(x,f(x))\,|\, x\in U\}\subseteq \mathbb{R}^{n+1}$，
  这是 $\mathbb{R}^{n+1}$ 的一个 $n$ 维嵌入子流形。其有全局参数化 $X:U\to \mathbb{R}^{n+1}$
  为 $X(u)=(u,f(u))$，这被称为\emph{图像参数化}，对应的 $M$ 上的坐标
  $(u^1,\dots,u^n)$ 被称为\emph{图像坐标}。在图像坐标中，$\Gamma(f)$
  的诱导度量为
  \[
    X^*\bar g=X^*\bigl(
      (dx^1)^2+\cdots+(dx^{n+1})^2
    \bigr)=(du^1)^2+\cdots+(du^n)^2+df^2.
  \]
  应用到 $\mathbb{S}^2$ 的上半球面，有参数化 $X: \mathbb{B}^2\to \mathbb{R}^3$
  为
  \[
    X(u,v)=\bigl(u,v,\sqrt{1-u^2-v^2}\bigr),
  \]
  于是 $\mathbb{S}^2$ 上的圆度量在图像坐标 $(u,v)$ 中表示为
  \begin{align*}
    \mathring{g}&=X^*\bar g=(du)^2+(dv)^2+d\bigl(\sqrt{1-u^2-v^2}\bigr)^2\\
    &=\frac{(1-v^2)\d u^2+(1-u^2)\d v^2+2uv\d u\d v}{1-u^2-v^2}.
  \end{align*}
\end{example}

\begin{example}[旋转曲面]\label{exa:surface of revolution}
  令 $H$ 是半平面 $\{(r,z)\,|\, r>0\}$，$C\subseteq H$ 是嵌入的 $1$-维子流形。定义
  由 $C$ 确定的\emph{旋转曲面}为子集 $S_C\subseteq \mathbb{R}^3$，
  \[
    S_C=\left\{
      (x,y,z)\middle| \left(
        \sqrt{x^2+y^2},z
      \right)\in C
    \right\}.
  \]
  集合 $C$ 被称为\emph{生成曲线}。每个关于 $C$ 的光滑局部参数化 $\gamma(t)=(a(t),b(t))$
  都导出了 $S_C$ 的一个光滑局部参数化
  \[
    X(t,\theta)=\bigl(a(t)\cos\theta,a(t)\sin\theta,b(t)\bigr),
  \]
  只要 $(t,\theta)$ 限制在某个足够小的开集中。此时 $t$-坐标曲线 $t\mapsto X(t,\theta_0)$
  被称为\emph{经线}，$\theta$-坐标曲线 $\theta\mapsto X(t_0,\theta)$ 被称为\emph{纬圈}。
  $S_C$ 上的诱导度量为
  \begin{align*}
    X^*\bar g&=d(a(t)\cos\theta)^2+d(a(t)\sin\theta)^2+d(b(t))^2\\
    &=\bigl(a'(t)\cos\theta\d t-a(t)\sin\theta\d\theta\bigr)^2+\bigl(a'(t)\sin\theta\d t+a(t)\cos\theta\d\theta\bigr)^2\\
    &\quad+(b'(t))^2\d t^2\\
    &=\bigl(a'(t)^2+b'(t)^2\bigr)\d t^2+a(t)^2\d\theta^2.
  \end{align*}
  特别地，如果 $\gamma$ 是\emph{单位速度曲线}(即对于任意 $t$ 都有 $|\gamma'(t)|^2=a'(t)^2+b'(t)^2=1$)，
  那么上述诱导度量为 $dt^2+a(t)^2\d\theta^2$。

  下面是一些具体的例子。
  \begin{itemize}[nosep]
    \item 如果 $C$ 是半圆 $r^2+z^2=1$，由 $\gamma(\varphi)=(\sin\varphi,\cos\varphi)\ (0<\varphi<\pi)$
    参数化，那么 $S_C$ 是单位球面(减去南北极点)。映射 $X(\varphi,\theta)=(\sin\varphi\cos\theta,\sin\varphi\cos\theta,\cos\varphi)$
    被称为\emph{球坐标参数化}，此时诱导度量为 $\d\varphi^2+\sin^2\varphi\d\theta^2$。
    \item 如果 $C$ 是圆 $(r-2)^2+z^2=1$，由 $\gamma(t)=(2+\cos t,\sin t)$ 参数化，那么我们
    得到环面，诱导度量是 $\d t^2+(2+\cos t)^2\d\theta^2$。
    \item 如果 $C$ 是 $\gamma(t)=(1,t)$ 参数化的竖直线，那么 $S_C$ 是单位圆柱面 $x^2+y^2=1$，
    此时诱导度量是 $\d t^2+\d\theta^2$。注意到，这意味着此时 $X: \mathbb{R}^2\to \mathbb{R}^3$
    是等距浸入。
  \end{itemize}
\end{example}

\begin{example}[$n$-环面]
  光滑覆盖映射 $X:\mathbb{R}^n\to \mathbb{T}^n$ 定义为 $X(u^1,\dots,u^n)=(e^{iu^1},\dots,e^{iu^n})$，
  那么 $X$ 可以限制到 $\mathbb{R}^n$ 的足够小的开子集上，成为一个光滑局部参数化，
  并且诱导度量在 $(u^i)$ 坐标中是 Euclid 度量，所以 $\mathbb{T}^n$ 上的诱导度量是平坦的。
\end{example}

\subsection{黎曼积}

如果 $(M_1,g_1)$ 和 $(M_2,g_2)$ 是黎曼流形，那么积流形 $M_1\times M_2$ 有一个
自然的黎曼度量 $g_1\oplus g_2$，被称为\emph{黎曼积}，定义为
\begin{equation}
  g_{(p_1,p_2)}\bigl((v_1,v_2),(w_1,w_2)\bigr)=g_{1}|_{p_1}(v_1,w_1)+g_{2}|_{p_2}(v_2,w_2).
\end{equation}
其中 $(v_1,v_2)$ 和 $(w_1,w_2)$ 是 $T_{p_1}M_1\oplus T_{p_2}M_2$ 的元素，
这自然地等同为 $T_{(p_1,p_2)}(M_1\times M_2)$ 的元素。$M_1$ 的光滑坐标卡
$(x^1,\dots,x^n)$ 和 $M_2$ 的光滑坐标卡 $(x^{n+1},\dots,x^{n+m})$
给出了 $M_1\times M_2$ 的光滑坐标卡 $(x^1,\dots,x^{n+m})$，在这个坐标卡中，黎曼积
$g_1\oplus g_2$ 的局部表示为 $g=g_{ij}\d x^i\d x^j$，其中 $(g_{ij})$ 是分块对角矩阵，前 $n$ 行 $n$ 列是 $g_1$ 的局部表示，后 $m$ 行 $m$ 列是 $g_2$ 的局部表示
\[
  (g_{ij})=
  \begin{pmatrix}
    (g_1)_{ab} & 0\\
    0 & (g_2)_{cd}
  \end{pmatrix}.
\]
其中 $a,b$ 从 $1$ 到 $n$，$c,d$ 从 $n+1$ 到 $n+m$。

积度量有一个重要的推广。假设 $(M_1,g_1)$ 和 $(M_2,g_2)$ 是黎曼流形，$f:M_1\to \mathbb{R}^+$ 是
严格正值的光滑函数。定义\emph{扭曲积} $M_1\times_f M_2$ 为积流形 $M_1\times M_2$ 配备
黎曼度量 $g=g_1\oplus f^2 g_2$，定义为
\[
  g_{(p_1,p_2)}\bigl((v_1,v_2),(w_1,w_2)\bigr)=g_{1}|_{p_1}(v_1,w_1)+f(p_1)^2 g_{2}|_{p_2}(v_2,w_2).
\]
其中 $(v_1,v_2),(w_1,w_2)\in T_{p_1}M_1\oplus T_{p_2}M_2$。
非常多类型的度量都可以这样构造，下面是几个简单的例子。

\begin{example}[扭曲积]
  \mbox{}
  \begin{enumerate}
    \item 当 $f\equiv 1$ 时，扭曲积 $M_1\times_f M_2$ 就是黎曼积 $M_1\times M_2$。
    \item 每个旋转曲面都可以表示为某个扭曲积。令 $H$ 是半平面 $\{(r,z)\,|\, r>0\}$，$C\subseteq H$ 是嵌入的 $1$-维子流形。定义
    $S_C\subseteq \mathbb{R}^3$ 是对应的旋转曲面。将 $C$ 配备 $H$ 上的诱导度量，
    $\mathbb{S}^1$ 配备标准度量。令 $f:C\to \mathbb{R}$ 是到 $z$-轴的距离
    $f(r,z)=r$。那么我们证明 $S_C$ 等距于扭曲积 $C\times_f \mathbb{S}^1$。
    \begin{proof}
      设 $\gamma(t)=(a(t),b(t))$ 是 $C$ 的一个光滑局部参数化，那么 $S_C$
      的一个光滑局部参数化为
      \[
        X(t,\theta)=(a(t)\cos\theta,a(t)\sin\theta,b(t)).
      \]
      由 \ref{exa:surface of revolution} 可知，$S_C$ 上的度量
      有表示
      \[
        X^*\bar g=\bigl(a'(t)^2+b'(t)^2\bigr)\d t^2+a(t)^2\d\theta^2.
      \]
      定义 $F: C\times_f \mathbb{S}^1\to S_C$ 为 
      \[
        F\bigl((r,z),e^{i\theta}\bigr)=\bigl(r\cos\theta,r\sin\theta,z\bigr).
      \]
      显然 $F$ 是一个微分同胚。显然 $(t,\theta)$ 构成 $C\times_f \mathbb{S}^1$ 的一个光滑坐标，
      在这个坐标中，$F$ 有坐标表示
      \[
        F(t,\theta)=\bigl(a(t)\cos\theta,a(t)\sin\theta,b(t)\bigr)=X(t,\theta).
      \]
      $C\times_f \mathbb{S}^1$ 上的度量有表示
      \[
        g=\bigl(a'(t)^2+b'(t)^2\bigr)\d t^2+a(t)^2\d\theta^2=X^*\bar g,
      \]
      这就表明 $F$ 是等距的。
    \end{proof}

    \item 令 $\rho$ 是 $\mathbb{R}^+\subseteq \mathbb{R}$ 上的标准坐标函数，
    那么 映射 $\Phi(\rho,\omega)=\rho\omega$ 给出了从扭曲积 $\mathbb{R}^+\times_\rho \mathbb{S}^{n-1}$
    到 $\mathbb{R}^n\setminus\{0\}$ 的等距。
  \end{enumerate}
\end{example}

\subsection{黎曼浸没}

与黎曼流形的子流形和积不同，黎曼流形的商只有在非常特殊的情况下才能继承一个度量。
本节我们来看什么时候出现这种情况。

假设 $M$ 和 $\wtilde M$ 是光滑流形，$\pi:\wtilde M\to M$ 是光滑浸没，
$\tilde g$ 是 $\wtilde M$ 上的黎曼度量。根据浸没水平集定理，每个纤维
$\wtilde M_y=\pi^{-1}(y)$ 都是 $\wtilde M$ 的一个恰当嵌入的子流形。对于每个
$x\in\wtilde M$，我们定义 $T_x\wtilde M$ 的两个子空间：
\emph{$x$ 处的竖直切空间}为
\[
  V_x=\ker \d\pi_x=T_x\bigl(\wtilde M_{\pi(x)}\bigr)
\]
(也即与 $x$ 的纤维相切的空间)，以及\emph{$x$ 处的水平切空间}为
\[
  H_x=V_x^\bot.
\]
于是 $T_x\wtilde M$ 有正交直和分解 $T_x\wtilde M=V_x\oplus H_x$。
注意到竖直空间对于每个浸没都可以定义，这与度量是无关的。但是水平空间依赖于度量。

$\wtilde M$ 上的向量场被称为\emph{水平向量场}，如果其在每个点处的值都在
水平空间中，类似地定义\emph{竖直向量场}。给定 $M$ 上的向量场 $X$，如果
存在 $\wtilde M$ 上的向量场 $\wtilde X$ 使得 $\wtilde X$ 是水平的，并且
和 $X$ 是 $\pi$-相关的 (即对于每个 $x\in\wtilde M$，有 $\d\pi_x(\wtilde X_x)=X_{\pi(x)}$)，
那么我们称 $\wtilde X$ 是 $X$ 的\emph{水平提升}。

下面的命题是在黎曼浸没上做计算的主要工具。

\begin{proposition}[水平向量场的性质]
  令 $\wtilde M$ 和 $M$ 是光滑流形，$\pi:\tilde M\to M$ 是光滑浸没，$\tilde g$
  是 $\wtilde M$ 上的黎曼度量。
  \begin{enumerate}
    \item $\wtilde M$ 上的每个光滑向量场 $W$ 都可以唯一表示为 $W=W^H+W^V$，
    其中 $W^H$ 是光滑水平的，$W^V$ 是光滑竖直的。
    \item $M$ 上的每个光滑向量场都有唯一的光滑水平提升。
    \item 对于每个 $x\in \wtilde M$ 和 $v\in H_x$，都存在一个向量场 $X\in \mathfrak X(M)$
    使得其水平提升 $\wtilde X$ 满足 $\wtilde X_x=v$。
  \end{enumerate}
\end{proposition}
\begin{proof}
  (1) 任取 $p\in \wtilde M$，因为 $\pi$ 是光滑浸没，根据秩定理，存在以 $p$ 为中心
  的光滑坐标卡 $(\wtilde U,(x^i))$ 和以 $\pi(p)$ 为中心的光滑坐标卡 $(U,(u^j))$，
  使得 $\pi$ 有坐标表示
  \[
    \pi\bigl(x^1,\dots,x^n,x^{n+1},\dots,x^m\bigr)=\bigl(
      x^1,\dots,x^n
    \bigr).
  \]
  对于任意 $q\in \wtilde U$，这表明竖直空间 $V_q$ 由坐标向量 $(\partial_{x^{n+1}}|_q,\dots,\partial_{x^m}|_q)$ 张成。
  对标架 $(\partial_{n+1},\dots,\partial_m,\partial_1,\dots,\partial_n)$
  使用 Gram-Schmidt 正交化，得到 $\wtilde U$ 上的光滑正交标架 $(E_1,\dots,E_m)$，
  并且对于每个 $q\in \wtilde U$，$V_q$ 由 $(E_1|_q,\dots,E_{m-n}|_q)$ 张成。
  因此 $H_q$ 由 $(E_{m-n+1}|_q,\dots,E_m|_q)$ 张成。

  对于任意光滑向量场 $W\in \mathfrak X(\wtilde M)$，在每个点 $q\in\wtilde M$
  处，$W_q$ 可以唯一表示为一个竖直向量和水平向量的和，这定义了一个分解 $W=W^V+W^H$。
  下面我们只需要说明 $W^V$ 和 $W^H$ 都是光滑的。在每个点附近，都可以构造上述
  的正交标架，使得 $W=W^1E_1+\dots+W^mE_m$，其中 $W^i$ 是光滑函数。于是
  $W^V=W^1E_1+\dots+W^{m-n}E_{m-n}$，$W^H=W^{m-n+1}E_{m-n+1}+\dots+W^mE_m$，
  这就表明 $W^V$ 和 $W^H$ 都是光滑的。

  (2) 任取 $X\in\mathfrak X(M)$，若 $\wtilde X$ 是 $X$ 的水平提升，那么对于每个
  $x\in \wtilde M$，有 $\wtilde X_x\in (\ker\d\pi_x)^\bot$ 并且
  $\d\pi_x(\wtilde X_x)=X_{\pi(x)}$。由于 $\d\pi_x$ 限制在 $H_x$ 上是一个线性同构，所以 $\wtilde X_x$ 是唯一确定的。
  于是水平提升 $\wtilde X$ 是唯一的。下面说明 $\wtilde X$ 是光滑的。
  任取 $p\in \wtilde M$，在 $p$ 的某个邻域 $\wtilde U$ 上构造上述的正交标架 $(E_1,\dots,E_m)$。
  对于任意 $q\in\wtilde U$，$\d\pi_q$ 限制在 $H_q$ 上是一个恒等映射，
  所以 $\wtilde X_q$ 可以表示为 $\wtilde X_q=X^1(\pi(q))E_{m-n+1}|_q+\cdots+X^n(\pi(q))E_m|_q$，
  这表明 $\wtilde X$ 在 $\wtilde U$ 上是光滑的。

  (3) 由于 $\d\pi_x$ 限制在 $H_x$ 上的同构，所以考虑 $v$ 的像 $\d\pi_x(v)\in T_{\pi(x)}M$，
  取一个光滑向量场 $X\in\mathfrak X(M)$，满足 $X_{\pi(x)}=\d\pi_x(v)$。
  令 $\wtilde X$ 为 $X$ 的水平提升，则 $\d\pi_x(\wtilde X_x)=X_{\pi(x)}=\d\pi_x(v)$
  且 $\wtilde X_x\in H_x$，所以 $\wtilde X_x=v$。
\end{proof}

上述命题的 (3) 表明 $\wtilde M$ 的某个点处的每个水平向量都可以延拓为一个水平提升，
这对于计算来说非常有用。但是需要注意，并不意味着 $\wtilde M$ 上的每个水平向量场都可以
成为一个水平提升。

现在我们可以把黎曼流形的某些商流形配备一个度量，使得商映射成为黎曼浸没。

假设 $\bigl(\wtilde M,\tilde g\bigr)$ 和 $(M,g)$ 是黎曼流形，$\pi:\wtilde M\to M$ 是一个光滑浸没。我们说 $\pi$ 是一个\emph{黎曼浸没}，
如果对于每个 $x\in\wtilde M$，微分 $\d\pi_x$ 都可以限制为等距同构 $H_x\to T_{\pi(x)}M$。
换句话说，当 $v,w\in H_x$ 的时候，总有 $\wtilde g_x(v,w)=g_{\pi(x)}(\d\pi_x(v),\d\pi_x(w))$。

给定一个黎曼流形 $\bigl(\wtilde M,\tilde g\bigr)$ 和满射的浸没 $\pi:\wtilde M\to M$，
实际上几乎不存在 $M$ 上的度量 $g$ 使得 $\pi$ 是一个黎曼浸没。这不难看出是为什么：
假设 $\pi$ 是一个黎曼浸没，如果 $p_1,p_2\in\wtilde M$ 处于同一个纤维 $\pi^{-1}(y)$
中，那么线性映射 $\bigl(\d\pi_{p_i}|_{H_{p_i}}\bigr)^{-1}:T_{y}M\to H_{p_i}$
需要把 $H_{p_i}$ 中的内积 $\tilde g_{p_i}$ 拉回到 $T_{y}M$ 上的同一个内积 $g_y$。

然而，在一个特殊的情况下，这样的度量是存在的。假设 $\pi:\wtilde M\to M$ 是一个满射光滑浸没，
并且 $G$ 是作用在 $\wtilde M$ 上的一个群。我们说这个作用是\emph{竖直的}，如果
每个 $\varphi\in G$ 把每个纤维送到自身，也即对于每个 $p\in\wtilde M$，
有 $\pi(\varphi\cdot p)=\pi(p)$。这个作用是\emph{在纤维上传递的}，
如果对于每个 $p,q\in\wtilde M$ 且 $\pi(p)=\pi(q)$，都存在 $\varphi\in G$ 使得 $\varphi\cdot p=q$。

如果 $\wtilde M$ 还配备一个黎曼度量，并且对于每个 $\varphi\in G$，映射
$x\mapsto \varphi\cdot x$ 是等距映射，那么这个作用被称为\emph{等距作用}，
并且这个度量被称为\emph{$G$-不变的}。在这种情况下，只要这个作用是有效的
(即 $G$ 中不同的元素定义了不同的等距)，那么我们就可以把 $G$ 视为 $\Iso(\wtilde M,g)$
的一个子群。

\begin{theorem}
  令 $\bigl(\wtilde M,\tilde g\bigr)$ 是黎曼流形，$\pi:\wtilde M\to M$
  是满射的光滑浸没，$G$ 是在 $\wtilde M$ 上的一个作用。如果这个作用是等距的、
  竖直的并且在纤维上是传递的，那么存在 $M$ 上唯一的黎曼度量 $g$ 使得 $\pi$ 是一个黎曼浸没。
\end{theorem}
\begin{proof}
  首先我们说明 $M$ 上的度量 $g$ 是存在且唯一的。任取 $y\in M$，对于每个 $p\in\pi^{-1}(y)$，
  $\d\pi_p$ 限制在水平空间 $H_p$ 上是一个线性同构，所以 $\bigl(\d\pi_p|_{H_p}\bigr)^{-1}:T_yM\to H_p$
  是一个线性同构。
  因此，如果存在一个度量 $g$ 使得 $\pi$ 是一个黎曼浸没，那么对于每个 $v,w\in T_yM$ 和任意 $p\in \pi^{-1}(y)$，
  必须有 
  \[
    g_y(v,w)=\tilde g_p\left(
      \bigl(\d\pi_p|_{H_p}\bigr)^{-1}(v),\bigl(\d\pi_p|_{H_p}\bigr)^{-1}(w)
    \right).
  \]
  接下来我们说明上式是良定义的，也即和 $\pi^{-1}(y)$ 中的点 $p$ 无关。
  
  任取 $\varphi\in G$，记等距 $F_\varphi:\wtilde M\to\wtilde M$ 为 
  $x\mapsto \varphi\cdot x$，由于 $\varphi$ 是竖直的，所以 $\pi\circ F_\varphi=\pi$。
  任取 $p,q\in \pi^{-1}(y)$，由于 $G$ 在纤维上传递，所以存在 $\varphi\in G$ 使得 $F_\varphi(p)=q$。
  这表明 $\d\pi_q\circ \d (F_\varphi)_p=\d\pi_p$，
  于是 $\bigl(\d\pi_q|_{H_q}\bigr)^{-1}=(\d (F_\varphi)_p|_{H_p})\circ \bigl(\d\pi_p|_{H_p}\bigr)^{-1}$。
  那么对于任意 $v,w\in T_yM$，有
  \begin{align*}
    &\quad \tilde g_{q}\left(
      \bigl(\d\pi_q|_{H_q}\bigr)^{-1}(v),\bigl(\d\pi_q|_{H_q}\bigr)^{-1}(w)
    \right)\\
    &=\tilde g_{q}\left(
      (\d (F_\varphi)_p|_{H_p})\circ \bigl(\d\pi_p|_{H_p}\bigr)^{-1}(v),(\d (F_\varphi)_p|_{H_p})\circ \bigl(\d\pi_p|_{H_p}\bigr)^{-1}(w)
    \right)\\
    &=\tilde g_p\left(
      \bigl(\d\pi_p|_{H_p}\bigr)^{-1}(v),\bigl(\d\pi_p|_{H_p}\bigr)^{-1}(w)
    \right),
  \end{align*}
  最后一个等式是因为 $F_\varphi$ 是等距映射。这就表明 $g$ 的定义是良定义的，且
  是唯一的。】

  不难验证 $g$ 是黎曼度量，且根据 $g$ 的构造，$\pi$ 是一个黎曼浸没。
\end{proof}

\begin{corollary}
  设 $\bigl(\wtilde M,\tilde g\bigr)$ 是黎曼流形，$G$ 是一个 Lie 群，
  $G$ 光滑、自由、恰当且等距地作用在 $\wtilde M$ 上。那么轨道空间 $M=\wtilde M/G$
  有一个唯一的光滑流形结构和黎曼度量，使得 $\pi$ 是黎曼浸没。
\end{corollary}

\subsection{黎曼覆盖}

\section{黎曼流形上的基本构造}

每个黎曼度量都可以导出流形上的非常丰富的结构，不仅仅是一些显然的向量
的长度和夹角的概念。在本节我们描述最基本的构造。本节中 $M$ 始终表示
一个带边或者无边光滑流形。

\subsection{指标上升和指标下降}

黎曼度量的一个重要且基本的作用是允许我们在向量和余向量之间转换。
给定 $M$ 上的一个黎曼度量 $g$，我们可以定义丛同态 $\hat g:TM\to T^*M$
为
\[
  \hat g(v)(w)=g_p(v,w),\quad \forall v,w\in T_pM, p\in M.
\]
如果 $X$ 和 $Y$ 是 $M$ 上的光滑向量场，那么上述 $\hat g$ 诱导出
向量丛截面之间的映射 $\hat g:\mathfrak X(M)\to \mathfrak X^*(M)$，
满足 
\[
  \hat g(X)(Y)=g(X,Y).
\]
由于 $\hat g(X)$ 在 $Y$-分量上是 $C^\infty(M)$ 线性的，所以 $\hat g(X)$
确实是 $M$ 上的一个光滑余向量场。并且 $\hat g$ 也是 $C^\infty(M)$ 线性的，
所以 $\hat g:TM\to T^*M$ 确实是一个丛同态。

给定光滑局部标架 $(E_i)$ 和其对偶标架 $(\varepsilon^i)$，设 $g$
的局部表示为 $g=g_{ij}\varepsilon^i\varepsilon^j$。如果 $X=X^iE_i$
是光滑向量场，那么余向量场 $\hat g(X)$ 有坐标表示
\[
  \hat g(X)=\bigl(g_{ij}X^i\bigr)\varepsilon^j.
\]
这表明 $\hat g$ 在任意局部标架中的矩阵表示为 $(g_{ij})$，这与 $g$ 的局部表示矩阵相同。

给定一个向量场 $X$，余向量场 $\hat g(X)$ 的分量通常记为
\[
  X_j=g_{ij}X^i,
\]
于是
\[
  \hat g(X)=X_j\varepsilon^j.
\]
我们说 $\hat g(X)$ 是 $X$ 通过\emph{指标下降}得到的，通常把 $\hat g(X)$ 记为 $X^\flat$。

因为矩阵 $(g_{ij})$ 在每个点处都是可逆的，所以映射 $\hat g$ 也是可逆的，
并且 $\hat g^{-1}$ 的矩阵就是 $(g_{ij})$ 的逆矩阵，我们记为 $(g^{ij})$。
于是 $g^{ij}g_{jk}=\delta^i_k$。$(g_{ij})$ 是对称阵表明 $(g^{ij})$ 也是对称阵。
那么在局部标架中，逆映射 $\hat g:T^*M\to TM$ 有坐标表示
\[
  \hat g^{-1}(\omega)=g^{ij}\omega_j E_i,
\]
一般，我们记 $\omega^i=g^{ij}\omega_j$，此时便有 $\hat g^{-1}(\omega)=\omega^i E_i$。
我们说 $\hat g^{-1}(\omega)$ 是 $\omega$ 通过\emph{指标上升}得到的，
通常把 $\hat g^{-1}(\omega)$ 记为 $\omega^\sharp$。这两个互逆的丛同构
$\flat$ 和 $\sharp$ 被称为\emph{音乐同构}。

$\sharp$ 算符最重要的应用是把梯度算符延拓到了黎曼流形上。如果 $g$ 是 $M$
上的黎曼度量，$f:M\to \mathbb{R}$ 是光滑函数，定义\emph{$f$ 的梯度}是
$\d f$ 的指标上升，即 $\grad f=(\d f)^\sharp$。利用这个定义，
有 $(\grad f)^\flat=\d f$，所以 $\grad f$ 可以刻画为
\begin{equation}
  \d f_p(w)=\left\langle \grad f|_p,w\right\rangle,\quad \forall w\in T_pM, p\in M.
\end{equation}
此外，$\grad f$ 有局部坐标表示
\[
  \grad f=(g^{ij}E_j f)E_i.
\]
因此，如果 $(E_i)$ 是正交标架，那么 $\grad f=\sum_i (E_i f) E_i$，
也即 $\grad f$ 的系数和 $\d f$ 的系数是一样的。

下面的命题表明，黎曼流形上的梯度和 Euclid 空间上的梯度有相同的几何解释。
如果 $f:M\to \mathbb{R}$ 是光滑实值函数，回顾一个点 $p\in M$ 被称为\emph{正则点}，
如果 $\d f_p\neq 0$，否则被称为\emph{临界点}。水平集 $f^{-1}(c)$ 被称为\emph{正则水平集}，
如果 $f^{-1}(c)$ 中的每个点都是正则点。

\begin{proposition}
  假设 $(M,g)$ 是黎曼流形，$f\in C^\infty(M)$，$\mathcal{R}\subseteq M$
  是 $f$ 的正则点的集合。对于每个 $c\in \mathbb{R}$，集合 $M_c=f^{-1}(c)\cap \mathcal{R}$
  如果非空，那么 $M_c$ 是 $M$ 中的嵌入光滑超曲面，并且 $\grad f$ 与 $M_c$ 处处正交。
\end{proposition}
\begin{proof}
  
\end{proof}










